

\definecolor{bg}{rgb}{0.95,0.95,0.95}
\definecolor{myBGcolor}{RGB}{204,204,255}


\definecolor{mygreen1}{RGB}{34, 170, 141}
\definecolor{mygreen2}{RGB}{37,185,111}
\definecolor{mygreen3}{RGB}{105,225,165}
\definecolor{mygreen4}{RGB}{43,213,177}
\definecolor{mygreen5}{RGB}{128,230,208}
\definecolor{mygreen6}{RGB}{185,241,229}
\definecolor{mygreen7}{RGB}{215,247,240}
\definecolor{mygreen8}{RGB}{88,214,58}
\definecolor{mygreen9}{RGB}{131,208,100}
\definecolor{mygreen10}{RGB}{204,255,199}
\definecolor{mygreen11}{RGB}{168,255,159}
\definecolor{mygreen12}{RGB}{218,255,218}


\definecolor{myred1}  {RGB}{219, 63, 61}
\definecolor{myred2} {RGB}{252, 118, 131}
\definecolor{myred3} {RGB}{250, 66, 84}
\definecolor{myred4} {RGB}{253, 177, 184}
\definecolor{myred5} {RGB}{244, 60, 82}
\definecolor{myred6} {RGB}{205, 63, 80}
\definecolor{myred7} {RGB}{253, 177, 179}
\definecolor{myred8} {RGB}{253, 136, 127}
\definecolor{myred9} {RGB}{255, 213, 213}


\definecolor{myyellow1}  {RGB}{251, 251, 139}
\definecolor{myyellow2}  {RGB}{226, 226, 8}
\definecolor{myyellow3} {RGB}{252, 252, 168}
\definecolor{myyellow4} {RGB}{253, 253, 203}
\definecolor{myyellow5} {RGB}{253, 253, 19}


\definecolor{myblue1} {RGB}{104, 203, 248}
\definecolor{myblue2} {RGB}{12, 168, 238}
\definecolor{myblue3} {RGB}{72, 192, 246}
\definecolor{myblue4} {RGB}{123, 210, 249}
\definecolor{myblue5} {RGB}{120, 214, 201}
\definecolor{myblue6} {RGB}{93, 241, 171}
\definecolor{myblue7} {RGB}{64, 68, 244}


\definecolor{mypurple1}{RGB}{204,204,255}
\definecolor{mypurple2}{RGB}{137,99,245}
\definecolor{mypurple3}{RGB}{121,78,244}
\definecolor{mypurple4}{RGB}{161,70,236}
\definecolor{mypurple5}{RGB}{140,29,231}
\definecolor{mypurple6}{RGB}{192,132,242}
\definecolor{mypurple7}{RGB}{204,204,255}


% 图占位：\figholder{图题}{标签}{图的内容说明}
\newcommand{\figholder}[3]{%
  \begin{figure}[H]
    \centering
    \fbox{\parbox{0.86\textwidth}{\centering\small\linespread{1.3}\selectfont
      \vspace{1.2em}\textbf{【图占位】}\par\vspace{0.4em}#3\par\vspace{1.2em}}}
    \caption{#1}
    \label{#2}
  \end{figure}}

% AI 输出节选：与 Prompt 的 quote 样式区分
\mdfdefinestyle{aioutputstyle}{%
  skipabove=\medskipamount, skipbelow=\medskipamount,
  leftline=true, rightline=false, topline=false, bottomline=false,
  linewidth=3pt, linecolor=blue!35, backgroundcolor=blue!3,
  innertopmargin=8pt, innerbottommargin=8pt,
  innerleftmargin=12pt, innerrightmargin=6pt,
}
\newenvironment{aioutput}{%
  \begin{mdframed}[style=aioutputstyle]%
  \small\linespread{1.3}\selectfont\textbf{AI 输出（节选）}\par\smallskip
}{%
  \end{mdframed}%
}

% 章首：本章学习目标（左侧色条）
\mdfdefinestyle{chapterguidestyle}{%
  skipabove=\bigskipamount, skipbelow=\bigskipamount,
  leftline=true, rightline=false, topline=false, bottomline=false,
  linewidth=4pt, linecolor=blue!55!black, backgroundcolor=blue!5,
  innertopmargin=10pt, innerbottommargin=10pt,
  innerleftmargin=14pt, innerrightmargin=10pt,
}
\newenvironment{chapterguide}{%
  \begin{mdframed}[style=chapterguidestyle]%
  {\large\bfseries\color{blue!55!black}本章学习目标}\par\smallskip
  \noindent 学完本章后，读者应当能够：\par
}{%
  \end{mdframed}%
}

\renewcommand{\lstlistingname}{代码清单}
\lstset{extendedchars=false,texcl=false,columns=fullflexible,keepspaces=true}

% 允许长代码路径在下划线处断行，避免溢出版心
\setlength{\emergencystretch}{3em}
\let\origunderscore\_
\renewcommand{\_}{\origunderscore\allowbreak}

% 概念图的统一样式（第 3 步新增插图使用）
% 注意：样式名不要用 out、id 等 TikZ 内置键名
\tikzset{
  pbox/.style={draw=black!55, rounded corners=2pt, fill=black!4, align=flush center,
    font=\small, inner sep=5pt, minimum height=9mm},
  hbox/.style={pbox, draw=blue!55!black, fill=blue!8},      % 人负责
  abox/.style={pbox, draw=orange!70!black, fill=orange!10}, % AI 负责
  gbox/.style={pbox, draw=green!45!black, fill=green!8},    % 通过 / 结果
  rbox/.style={pbox, draw=red!55!black, fill=red!7},        % 风险 / 错误
  tag/.style={font=\scriptsize, align=flush center, text=black!70},
  flow/.style={-Stealth, thick, draw=black!65},
  backflow/.style={-Stealth, dashed, draw=black!45},
}

% 左对齐的定宽列：避免中文在窄列中被两端对齐拉开
\newcolumntype{L}[1]{>{\raggedright\arraybackslash}p{#1}}

% 插入 TikZ 图：超过版心宽度时按比例缩小，否则保持原尺寸
\newsavebox{\fitfigbox}
\newcommand{\fitinput}[1]{%
  \sbox{\fitfigbox}{\input{#1}}%
  \ifdim\wd\fitfigbox>\linewidth
    \resizebox{\linewidth}{!}{\usebox{\fitfigbox}}%
  \else
    \usebox{\fitfigbox}%
  \fi}
